\begin{proposition}[Hard-prior loss of the certified simulator]\label{prop:step-013-simulator-hardness}
Under Proposition~\ref{prop:step-004-finite-hard-prior} and
Proposition~\ref{prop:step-012-source-membership}, if the total
simulator \(B=B_{\mu_{N,M},A}\) is instantiated at their identical
admissible \(N,M\) and public prior \(\mu_{N,M}\), then

\[
\mathbb E_{\substack{(T,Q)\sim\mu_{N,M}\\
                      Z\sim(Q^{\tau_T})^M\\ B}}
R_Q(B(Z),\tau_T)>\eta=2^{-8}.
\tag{9}
\]

The expectation includes the prior draw, exactly \(M\) i.i.d. realizable
one-block input rows conditional on that draw, all simulator preprocessing,
and any internal coins of \(A\) used by the simulator.
\end{proposition}

\begin{proof}
Proposition~\ref{prop:step-012-source-membership} gives every
membership clause needed by the universal learner quantifier in Proposition~\ref{prop:step-004-finite-hard-prior}: \(B\) is total on the
entire arbitrary labeled input space of exact size \(M\), its output is an
arbitrary member of \(\{0,1\}^{[N]}\), and it satisfies every event-level
one-row replacement inequality at privacy parameters

\[
\left(0.1,\frac{d_*}{M^2\log M}\right).
\tag{10}
\]

Both propositions use the same finite ordered domain \([N]\),
thresholds \(\tau_t\), population-risk convention, source constants, and
integers \(N,M\). Thus no convention, object, sample-size, or privacy bridge
remains to prove.

The prior in Proposition~\ref{prop:step-004-finite-hard-prior} was selected
before the universal learner
quantifier. Although the code of \(B\) uses the public
\(\mu_{N,M}\), it is therefore one of the prior-aware kernels already
covered by that quantifier. Substituting this exact \(B\) into (5) yields
(9) directly. The substitution neither conditions on overflow nor changes
the simulator on any input, so it introduces no residual. \end{proof}

\begin{proposition}[Prior-averaged product-risk transfer]\label{prop:step-013-product-lower-bound}
Under Proposition \ref{prop:step-013-simulator-hardness}, Propositions
\ref{prop:step-008-selected-risk-identity},
\ref{prop:step-009-overflow}, and
\ref{prop:step-011-risk-transfer}, the ideal experiment (3) obeys the exact
product-risk lower bound (4).
\end{proposition}

\begin{proof}
Define only for this proof

\[
\mathcal L_{\mathrm{act}}
:=
\mathbb E_{\substack{(T,Q)\sim\mu_{N,M}\\
                      Z\sim(Q^{\tau_T})^M\\ B}}
R_Q(B(Z),\tau_T)
\tag{11}
\]

and

\[
\mathcal L_{\mathrm{id}}
:=
\mathbb E_{\mathrm{ideal}}
R_{Q_J}(D_JH^{\mathrm{id}},\tau_{T_J}).
\tag{12}
\]

These are proof-local names for the two exact quantities already typed by
the dependencies. Proposition~\ref{prop:step-013-simulator-hardness}
gives

\[
\mathcal L_{\mathrm{act}}>2^{-8}.
\tag{13}
\]

Proposition~\ref{prop:step-011-risk-transfer} and the notation
\(p_{\mathrm{ov}}=\Pr(U>M)\) from (7) give

\[
\mathcal L_{\mathrm{id}}
\ge \mathcal L_{\mathrm{act}}-p_{\mathrm{ov}}.
\tag{14}
\]

Combining (13) with the weak inequality (14) preserves strictness:

\[
\mathcal L_{\mathrm{id}}
>2^{-8}-p_{\mathrm{ov}}.
\tag{15}
\]

Proposition~\ref{prop:step-009-overflow} gives the strict upper
bound \(p_{\mathrm{ov}}<2^{-9}\). Since subtraction reverses the comparison
in the subtracted term,

\[
2^{-8}-p_{\mathrm{ov}}
>2^{-8}-2^{-9}
=2\cdot2^{-9}-2^{-9}
=2^{-9}.
\tag{16}
\]

Equations (15)--(16) therefore yield

\[
\mathcal L_{\mathrm{id}}>2^{-9}.
\tag{17}
\]

The overflow event has been charged exactly once: it occurs only in (14),
and no later equality or bound adds another copy.

Finally, Proposition~\ref{prop:step-008-selected-risk-identity}
identifies (12), with no residual, as

\[
\mathcal L_{\mathrm{id}}
=
\mathbb E_{\substack{\boldsymbol\Xi\sim\mu_{N,M}^{\otimes k}\\
                      S^{\mathrm{id}}\sim
                      (P_{\boldsymbol Q}^{c_{\boldsymbol T}})^n\\
                      H^{\mathrm{id}}\sim A(S^{\mathrm{id}})}}
R_{P_{\boldsymbol Q}}(H^{\mathrm{id}},c_{\boldsymbol T}).
\tag{18}
\]

Substituting (18) into (17) proves (4). The equality in (18) is taken under
the unconditional ideal experiment; the conditional i.i.d. sample law is
invoked only after fixing \(\boldsymbol\Xi\), and no conditioning on
\(U\le M\) appears. \end{proof}
